\section{逐条指令 trace：不省略循环中的状态}

前面的四轮表只保留了每轮结果。下面把任务从入口到返回的 23 条动态指令全部列出。为了
控制表格宽度，地址省略共同的前导 \code{0x00}；“执行后”表示该指令成功提交后的体系结构
状态，未列出的寄存器和内存保持不变。

\begin{longtable}{|L{0.08\linewidth}|L{0.19\linewidth}|L{0.29\linewidth}|L{0.30\linewidth}|}
\hline
\rowcolor{BookTableHead}
序号 & 指令地址 & 执行动作 & 提交后的关键状态 \\ \hline
1 & \code{100000} & \code{MOVI r4,0} & \code{r4=0, PC=100004} \\ \hline
2 & \code{100004} & 从 \code{102000} 装入 7 & \code{r5=7, PC=100008} \\ \hline
3 & \code{100008} & \(0+7\) & \code{r4=7, PC=10000c} \\ \hline
4 & \code{10000c} & 输入指针加 4 & \code{r0=102004} \\ \hline
5 & \code{100010} & 计数减 1 & \code{r1=3} \\ \hline
6 & \code{100014} & \code{r1!=0}，采用分支 & \code{PC=100004} \\ \hline
7 & \code{100004} & 从 \code{102004} 装入 \(-2\) & \code{r5=-2} \\ \hline
8 & \code{100008} & \(7+(-2)\) & \code{r4=5} \\ \hline
9 & \code{10000c} & 输入指针加 4 & \code{r0=102008} \\ \hline
10 & \code{100010} & 计数减 1 & \code{r1=2} \\ \hline
11 & \code{100014} & \code{r1!=0}，采用分支 & \code{PC=100004} \\ \hline
12 & \code{100004} & 从 \code{102008} 装入 11 & \code{r5=11} \\ \hline
13 & \code{100008} & \(5+11\) & \code{r4=16} \\ \hline
14 & \code{10000c} & 输入指针加 4 & \code{r0=10200c} \\ \hline
15 & \code{100010} & 计数减 1 & \code{r1=1} \\ \hline
16 & \code{100014} & \code{r1!=0}，采用分支 & \code{PC=100004} \\ \hline
17 & \code{100004} & 从 \code{10200c} 装入 5 & \code{r5=5} \\ \hline
18 & \code{100008} & \(16+5\) & \code{r4=21} \\ \hline
19 & \code{10000c} & 输入指针加 4 & \code{r0=102010} \\ \hline
20 & \code{100010} & 计数减 1 & \code{r1=0} \\ \hline
21 & \code{100014} & \code{r1==0}，不分支 & \code{PC=100018} \\ \hline
22 & \code{100018} & 向 \code{102020} 存储 21 & \code{MEM32[102020]=21} \\ \hline
23 & \code{10001c} & 读取栈顶并返回 & \code{PC=000300, SP=1f0000} \\ \hline
\end{longtable}

动态指令数可以由结构独立复算。初始化执行一次；循环体包含
装入、加法、指针更新、计数更新和分支，共 5 条，执行四轮得到 20 条；最后存储和返回各
一次，所以总数为 \(1+20+2=23\)。算式与逐行编号一致，说明循环次数和退出路径没有遗漏。

这类自我校验正是使用具体程序的价值。抽象文字很难暴露“循环到底有几条动态指令”这样的
小错误，而带编号的 trace 会立即产生矛盾。

\subsection{事务数与指令数不是同一个数量}

每条指令至少需要一次取指事务，共 23 笔。四条动态 \code{LOAD} 再产生 4 笔数据读，
\code{STORE} 产生 1 笔数据写，\code{RET} 产生 1 笔栈读。因此不考虑启动前的人工准备，任务
总共发出：
\[
23+4+1+1=29
\]
笔外部事务。

算术指令也会读取寄存器，但寄存器堆位于处理器内部，不经过地址互连。把所有“读”都称为
内存事务会高估外部访问；反过来只数显式 \code{LOAD}/\code{STORE}，又会漏掉取指和返回
栈读取。

\subsection{若 RAM 响应时间不同，程序结果是否改变}

设启动存储读取需 2 个周期，RAM 读取需 3 个周期，RAM 写入需 2 个周期。教学处理器在
外部事务期间等待，因此总运行周期会增加，但只要每笔事务最终按同一顺序成功，23 条动态
指令提交后的寄存器和内存结果不变。

这就是功能抽象与性能事实的分界：软件依赖读取返回相应地址的字节，不应依赖每笔读取恰好
几个周期；若程序需要计时，必须另有可读计时器契约。

\subsection{三组反例：改变一个初值会发生什么}

\subsection{元素个数为零}

当前代码先执行一次 \code{LOAD}，然后才减少计数并判断。因此若启动时
\code{r1=0}，它仍会读取一个元素，随后计数变为 \code{0xffffffff}，分支继续，最终很可能
越过数组。这说明启动契约必须要求 \code{r1>0}，或者任务代码要在第一次装入前增加零长度
检查。

硬件不会替程序选择哪一种语义。若任务格式允许空数组，正确代码应先判断计数；若格式规定
至少一个元素，启动程序应验证输入。约束放在哪里，会决定错误由哪个层次报告。

\subsection{结果地址与输入重叠}

若启动程序把 \code{r2} 错设为 \code{0x00102008}，最终存储会覆盖第三个输入整数 11。由于
写发生在所有读取之后，本次结果仍是 21；但任务返回后输入数据已改变。若算法在写结果后
还要重读数组，行为就会变化。

“这一次碰巧正确”不能证明布局合法。装入契约若要求输入保持不变，就应让结果槽与输入范围
不重叠，并在建立参数时检查。

\subsection{数组末尾越过 RAM}

假设 \code{r0=0x001ffffc}、\code{r1=2}。第一轮能读取 RAM 最后 4 字节；指针随后变为
\code{0x00200000}，第二轮地址不在任何目标范围，互连返回 \code{UNMAPPED}。检查数组范围
时应验证：
\[
\texttt{count}\le
\frac{\texttt{RAM\_END}-\texttt{array\_base}}{4},
\]
并在乘法前避免 \code{count*4} 溢出。

\subsection{把器件承诺重新放回整机}

现在所有交互已经走读完成，可以进行一次横向复核：

\begin{longtable}{|L{0.17\linewidth}|L{0.25\linewidth}|L{0.25\linewidth}|L{0.23\linewidth}|}
\hline
\rowcolor{BookTableHead}
器件 & 保存或处理的事实 & 程序可依赖的接口 & 当前整机得到的承诺 \\ \hline
时钟与复位 & 边沿、复位有效状态 & 固定复位 PC & 确定起点与状态提交节拍 \\ \hline
处理器 & PC、SP、寄存器、译码状态 & 指令集与调用约定 & 可顺序推进的执行流 \\ \hline
启动存储 & 非易失字节阵列 & 只读物理范围 & 复位后立即存在的固件字节 \\ \hline
RAM & 易失可写单元 & 可读写物理范围 & 运行期可变字节空间 \\ \hline
地址互连 & 地址比较和响应路由 & 物理地址表与错误码 & 地址唯一选择目标 \\ \hline
\end{longtable}

最右列刻意使用有限承诺。RAM 不承诺断电保存，互连不承诺权限，处理器也不承诺当前任务
最终返回。准确列出这些限制，是为了让下一次增加硬件或程序时能够指出它究竟补上了哪一项
缺口。

\subsection{章末自检：能否独立重建这次运行}

读完本章后，应当能够不借助“系统自动完成”这样的句子回答以下问题。

\begin{enumerate}
\item 复位释放后，为什么第一笔取指访问启动存储而不是任务 RAM？
\item \code{LOAD r5,[r0+0]} 的指令地址、数据地址和目标寄存器号分别是什么？
\item 为什么最后一个人工写入字节完成不等于任务已经可以执行？
\item 哪些字段共同证明入口地址合法？哪个算术检查必须防止 32 位回绕？
\item 启动程序为什么要分别建立任务 SP 和自己的 SP？
\item PC 在什么确切时刻从启动程序范围进入任务范围？此前谁控制取指次序？
\item 结果写入 RAM、任务返回、操作者稳定读回分别由什么事件完成？
\item 如果任务进入无限循环，当前机器中的哪个器件能够强制改变它的 PC？
\end{enumerate}

前七问都能沿本章给出的地址、字段和事务找到具体答案。第八问的答案是“没有”。普通时钟
只推动当前指令流，当前也没有能够请求控制转移的外部部件。这一空缺不是概念定义，而是
机器在无限循环反例中的可观察事实。

\begin{problemframe}
\textbf{第一章得到的机器。} 它能从确定复位入口启动，验证一项由操作者预先写入 RAM 的
程序，为该程序建立参数和栈，逐条执行到结果写回，并在程序合作时返回启动代码。它尚不能
自动接收程序。下一章只增加一条外部字节通道，由“人工准备为何不可持续”推导接收、校验
与交接程序。
\end{problemframe}
